% Generated draft: requires mathematical and source review.
\documentclass[11pt]{article}
\usepackage[margin=25mm]{geometry}
\usepackage{fontspec}
\setmainfont{Latin Modern Roman}
\usepackage{amsmath,amssymb,mathtools,unicode-math}
\setmathfont{Latin Modern Math}
\usepackage{enumitem,needspace,xurl,hyperref}
\setlength{\parindent}{0pt}
\setlength{\parskip}{5pt}
\setlength{\emergencystretch}{3em}
\newcommand{\N}{\mathbb{N}}
\newcommand{\Z}{\mathbb{Z}}
\newcommand{\Q}{\mathbb{Q}}
\newcommand{\R}{\mathbb{R}}
\newcommand{\C}{\mathbb{C}}
\newcommand{\F}{\mathbb{F}}
\newcommand{\A}{\mathbb{A}}
\newcommand{\PP}{\mathbb{P}}
\newcommand{\E}{\mathbb{E}}
\newcommand{\eps}{\varepsilon}
\newcommand{\dd}{\,\mathrm{d}}
\begin{document}
\section*{UnsolvedMath ID: 700633}
\textbf{Source problem code:} catalogue entry
\subsection{Definitions and mathematical statement}
For a group \(G\), let \(G'=\langle x^{-1}y^{-1}xy:x,y\in G\rangle\) be its derived subgroup. Write \(x^G=\{g^{-1}xg:g\in G\}\). A group \(G\) is an \(FC\)-group if \(x^G\) is finite for every \(x\in G\). Let \(\mathcal F\) denote the class of all \(FC\)-groups. For a prime \(p\), the quasicyclic group is
\[
C_{p^\infty}=\{z\in\C^\times:z^{p^n}=1\text{ for some integer }n\geq 0\},
\]
with multiplication. Define
\[
\mathcal A=\{G:|G'|<\infty\},\qquad
\mathcal B=\{G\in\mathcal F:G'\cong C_{p^\infty}\text{ for some prime }p\}.
\]
The source calls a class of groups a \emph{direct variety} if it is closed under subgroups, factor groups, and direct products. The permitted direct products are not specified. To expose this unresolved convention, let \(\mathscr I\) be a collection of allowed index sets, and let \(\mathcal D_{\mathscr I}(\mathcal C)\) be the smallest isomorphism-closed class containing \(\mathcal C\) and closed under subgroups, quotients by normal subgroups, and products \(\prod_{i\in I}G_i\) whenever \(I\in\mathscr I\) and all factors belong to the class.

The question of L. A. Kurdachenko is whether, for the product convention intended by the source,
\[
\mathcal D_{\mathscr I}(\mathcal A\cup\mathcal B)=\mathcal F.
\]
This is a conditional formulation: the intended \(\mathscr I\) remains unspecified. The source asserts that \(\mathcal F\) is a direct variety, but \(FC\)-groups are not closed under unrestricted Cartesian products. They are closed under finite direct products. Consequently, the unrestricted interpretation is incompatible with the supplied closure assertion; the intended restriction must be checked before this can be treated as an unambiguous problem statement.

As background, the source attributes to P. Hall the assertion that finite groups together with abelian groups do not generate \(\mathcal F\) as a direct variety, and to L. A. Kurdachenko the assertion that \(\mathcal A\) alone does not generate \(\mathcal F\). These attributions retain the same unresolved product convention.
\subsection{Short English statement}
Kurdachenko asks whether every \(FC\)-group belongs to the direct variety generated by groups with finite derived subgroup together with \(FC\)-groups whose derived subgroup is quasicyclic. The meaning of the permitted direct products requires clarification: unrestricted Cartesian products conflict with the source's assertion that \(FC\)-groups form a direct variety. The supplied registry labels the problem open; that label has not been independently verified.
\subsection{Sources}
\begin{itemize}
\item Local UnsolvedMath source registry, permanent ID \(700633\), exact source problem code \texttt{catalogue entry}, title \emph{Kourovka Notebook Problem 10.16}. The supplied record attributes the problem to L. A. Kurdachenko, proposed in 1986, and labels it open.
\item \emph{The Kourovka Notebook (Unsolved Problems in Group Theory)}, problem 10.16, 10th issue (1986); supplied background also identifies the 21st issue edition (2026), arXiv identifier \texttt{1401.0300v48}. The supplied set description names E. I. Khukhro and V. D. Mazurov as editors.
\item Yu. M. Gorchakov, \emph{Groups with Finite Classes of Conjugate Elements}, Moscow, Nauka, 1978 (Russian). Cited by the source for the definition of a direct variety.
\item P. Hall, \emph{J. London Math. Soc.}, volume 34, no. 3 (1959), pp. 289--304. Cited by the source for the assertion concerning generation by finite and abelian groups; no article title is supplied.
\item L. A. Kurdachenko, \emph{Ukrain. Math. J.}, volume 39, no. 3 (1987), pp. 255--259. Cited by the source for the assertion concerning groups with finite derived subgroup; no article title is supplied.
\end{itemize}
\end{document}
